\documentclass[11pt,a4paper]{article}
\usepackage[margin=25mm]{geometry}
\usepackage[T1]{fontenc}
\usepackage[utf8]{inputenc}
\usepackage{lmodern,microtype,amsmath,amssymb,amsthm,booktabs,longtable,enumitem}
\usepackage[round,sort]{natbib}
\usepackage{xurl}
\usepackage[colorlinks=true,allcolors=blue]{hyperref}
\setlength{\emergencystretch}{2em}
\numberwithin{equation}{section}
\newtheorem{theorem}{Theorem}[section]
\newtheorem{proposition}[theorem]{Proposition}
\newtheorem{lemma}[theorem]{Lemma}
\newtheorem{corollary}[theorem]{Corollary}
\newtheorem{conjecture}[theorem]{Conjecture}
\newtheorem{maintheorem}[theorem]{Theorem}
\theoremstyle{definition}
\newtheorem{definition}[theorem]{Definition}
\newtheorem{fact}{Standard fact}
\renewcommand{\thefact}{I\arabic{fact}}
\theoremstyle{remark}
\newtheorem{remark}[theorem]{Remark}
\newcommand{\Fp}{\mathbb F_p}
\newcommand{\Fq}{\mathbb F_{p^2}}
\newcommand{\Fb}{\overline{\mathbb F}_p}
\newcommand{\Z}{\mathbb Z}
\newcommand{\Q}{\mathbb Q}
\newcommand{\Res}{\operatorname{Res}}
\newcommand{\Aut}{\operatorname{Aut}}
\newcommand{\Stab}{\operatorname{Stab}}
\newcommand{\den}{\operatorname{den}_{\mathrm{red}}}
\newcommand{\leg}[2]{\left(\frac{#1}{#2}\right)}
\newcommand{\Lp}[1]{\left(\frac{#1}{p}\right)}

\title{Supersingular poles and signed residues\\of prime-level Fricke Hauptmoduln}
\author{David Jordan}
\date{4 November 2026}
\begin{document}
\maketitle
\begin{abstract}
For $N\in\{2,3,5,7,13\}$, let $T_N=q^{-1}+O(q)$ be the normalized Hauptmodul for
$\Gamma_0(N)^+$. For every prime $p\ge5$ with $p\nmid N$, we determine the poles and
residues of the rational function representing $T_N|U_p$ modulo $p$. Its poles are
exactly the supersingular values with effective ramification index one. If $t_i$ is
such a value, its residue multiplied by
$\prod_{k\ne i}(t_i-t_k)^{(p+1)/e_k}$ equals $-\left(\frac{-N}{p}\right)$, where the
product includes all distinct supersingular values and $e_k$ is the corresponding
ramification index. The proof combines a Kaneko--Zagier derivative argument with
Baker's supersingular congruence and an eta-character calculation. The latter gives
$\bigl(\Delta(\tau)/\Delta(\ell\tau)\bigr)^{(p^2-1)/24}=\left(\frac{\ell}{p}\right)$
at every supersingular point of $X_0(\ell)$ for every prime $\ell\ne p$.
The ramification indices also admit an interpretation as node thicknesses in the
Deligne--Rapoport Fricke quotient. Exact computations provide finite certificates
for all admissible pairs with $p\le211$.
\end{abstract}
\noindent\textbf{Keywords:} supersingular elliptic curves; Fricke groups; Hauptmoduln; Hecke operators; eta quotients; congruences.

\section{Introduction}\label{sec:intro}
The reduction of $U_p$ applied to a modular function can encode the supersingular
locus as a divisor of poles. At level one, this relation is developed by \citet{dS94} and \citet[\S10]{KZ98}; \citet{Baker98} proves the
supersingular congruence that determines the sign in the level-one residue formula.
\citet{BL18} obtain corresponding formulas for Legendre modular
polynomials. We establish a uniform signed formula on the five prime-level Fricke
curves for which $X_0(N)$ itself has genus zero.

The residues contain information beyond the support of the supersingular divisor.
Since the reduced function vanishes at the cusp and has only simple poles, its
residues determine the entire function. An explicit residue formula therefore
reconstructs the reduced $U_p$ image from supersingular values and their weights,
and detects vanishing when no contributing value remains. The level-one theory
provides the model for this reconstruction; the Fricke quotient introduces
additional ramification, which must enter both the pole criterion and the product.

The distinction between ramified and unramified supersingular values is essential.
A ramified value is removed from the pole divisor, but remains in the residue
product with an exponent determined by its ramification index. The algebraic proof
makes this distinction through the vanishing order of a quotient involving the Hasse
invariant and the derivative of the Hauptmodul. Section~\ref{sec:geometry} explains
the same indices as thicknesses of nodes in an integral modular model.

The sign calculation also has a formulation independent of genus zero.
\citet[Theorem~1.2 and Remark~4]{Ono25} proves that the supersingular value of
$\Delta(\tau)/\Delta(\ell\tau)$ raised to $(p^2-1)/12$ is one, and exhibits
examples in which the half-exponent $(p^2-1)/24$ gives $-1$.
Theorem~\ref{thm:C} determines this remaining sign: it is
$\left(\frac{\ell}{p}\right)$ for every prime $\ell\ne p$, at every
supersingular point of $X_0(\ell)$. Thus the eta/Frobenius identity applies
without a genus-zero hypothesis, whereas its use in the explicit residue
formula requires the coordinates fixed below.

\subsection{Notation and normalization}
Throughout the pole and residue results,
$N\in\mathcal N:=\{2,3,5,7,13\}$ and $p\ge5$ is prime with $p\ne N$. Put
\[
a_N=\frac{24}{N-1},\qquad b_N=N^{12/(N-1)},\qquad
u_N=\left(\frac{\eta(\tau)}{\eta(N\tau)}\right)^{a_N},\qquad
T_N=u_N+a_N+b_N/u_N.
\]
The pairs $(a_N,b_N)$ are $(24,4096),(12,729),(6,125),(4,49),(2,13)$.
The Fricke involution sends $u_N$ to $b_N/u_N$, and
$u_N=q^{-1}-a_N+O(q)$, so $T_N=q^{-1}+O(q)$ has zero constant term.
These are the McKay--Thompson series of classes $2A,3A,5A,7A,13A$
\citep{CN79}; the proofs below use their modular properties.

We write $T=T_N$, $u=u_N$, $q=e^{2\pi i\tau}$, $\theta=q\,d/dq$,
$k=(p^2-1)/24$, $A=\overline E_{p-1}$ and $B=\overline E_{p+1}$.
Thus $A$ is the Hasse invariant, with $q$-expansion $1$, and $B$ has the
$q$-expansion of $E_2$ modulo $p$. Put
$\Delta=\eta^{24}$ and $\Delta_N=\Delta(\tau)\Delta(N\tau)$.
For $f=\sum c_nq^n$, define $f|U_p=\sum c_{pn}q^n$.
A bar denotes reduction modulo $p$. The finite fields are
$\Fp$, $\Fq$ and $\Fb$, and $\left(\frac{a}{p}\right)$ is the Legendre symbol.
When a modular form is evaluated at an elliptic curve, an invariant differential
is understood; all weight-zero ratios are independent of its choice.

The functions $u$ and $T$ remain coordinates on $X_0(N)_{\Fp}$ and
$X_0(N)^+_{\Fp}$ respectively. Indeed, the eta quotient is a modular unit with
divisor the difference of the two cusp sections, with unit leading coefficients
at both cusps away from $N$. It therefore gives the integral genus-zero coordinate;
the quotient by $u\mapsto b_N/u$ has coordinate $T$ and a single cusp.
This also fixes the normalization used in every product below.
\subsection{Weights}
Let $t_1,\dots,t_s\in\Fb$ be the distinct supersingular values of $T$. For a point
$x=(E,C)$ of $X_0(N)_{\Fb}$ over $t_i$ put
\[
m_x=\bigl|\Aut(E,C)/\{\pm1\}\bigr|,\qquad e_i=m_x\,\bigl|\Stab_{\langle w_N\rangle}(x)\bigr|.
\]
Call $t_i$ \emph{thin} if $e_i=1$ and put $D_{\rm thin}(T)=\prod_{e_i=1}(T-t_i)$. For thin $t_i$ let
\[
r_i=\Res_{T=t_i}R_{N,p}(T),\qquad W_i=\prod_{k\ne i}(t_i-t_k)^{(p+1)/e_k}.
\]
Section~\ref{sec:weights} shows that $e_i$ does not depend on the representative of the Fricke
orbit, that it is the effective (tame) inertia of the quotient map at $x$ (not the order of a group
containing $\pm1$), and that every exponent $(p+1)/e_k$ is an integer.

\subsection{Main results}
\begin{maintheorem}\label{thm:main}
Let $N\in\{2,3,5,7,13\}$ and let $p\ge5$ be a prime with $p\ne N$.
\begin{enumerate}[label=\textup{(\alph*)}]
\item \textbf{(Rationality.)} There is a unique $R_{N,p}\in\Fp(T)$ with $\overline{T|U_p}=R_{N,p}(\overline T)$. It vanishes at $T=\infty$.
\item \textbf{(Pole support.)} $R_{N,p}=P/D_{\rm thin}$ with $P\in\Fp[T]$,
$\deg P<\deg D_{\rm thin}$ and $\gcd(P,D_{\rm thin})=1$. Every pole is simple and has a nonzero residue. In particular $\den R_{N,p}=D_{\rm thin}$.
\item \textbf{(Residue as a Hasse derivative.)} If $t_i$ is thin and $x$ lies over $t_i$, then
\[
r_i=\frac{144\,(\theta T)(x)^{p+1}}{B(x)^2},\qquad \theta=q\frac{d}{dq},
\]
a ratio of two mod-$p$ modular forms of weight $2(p+1)$, evaluated at $x$.
\item \textbf{(Product.)} For thin $t_i$, with $x=(E,C)$ over $t_i$ and $E'=E/C$,
\[
r_iW_i=\frac{B^{p-1}}{\Delta_N^{\,k}}(x)=\Lp{N}\,c(E)\,c(E'),\qquad
c(E):=\frac{B(E,\omega)^{(p-1)/2}}{\Delta(E,\omega)^{k}}.
\]
\item \textbf{(Signed weighted residue formula.)} For every thin $t_i$,
\[
\boxed{\;r_iW_i=-\Lp{-N}\;}\qquad\text{equivalently}\qquad
r_i=-\Lp{-N}\prod_{k\ne i}(t_i-t_k)^{-(p+1)/e_k}.
\]
In particular $r_i\in\Fp^\times$ and $W_i\in\Fp^\times$, also when $t_i\notin\Fp$.
\item \textbf{(Hecke form.)} With $\Psi_{N,p}(X,Y)=(X^p-Y)(X-Y^p)+pH_{N,p}(X,Y)$ as in
Proposition~\ref{prop:Psi}, $\overline H_{N,p}(t^p,t)=\Lp{-N}W_i^{-1}$ if $t=t_i$ is thin, and
$\overline H_{N,p}(t^p,t)=0$ for every other $t\in\Fq$.
\end{enumerate}
\end{maintheorem}

\begin{corollary}[reconstruction]\label{cor:reconstruction}
The supersingular values and their indices determine the entire reduced function:
\begin{equation}\label{eq:partial-fractions}
R_{N,p}(T)=-\Lp{-N}\sum_{e_i=1}\frac{1}{W_i(T-t_i)}.
\end{equation}
Equivalently, if $D_{\rm thin}=\prod_f f$ is its factorization into monic
irreducibles over $\Fp$, then
\begin{equation}\label{eq:factor-reconstruction}
R_{N,p}(T)=\sum_{f\mid D_{\rm thin}} c_f\frac{f'(T)}{f(T)},\qquad
c_f=-\Lp{-N}W_i^{-1}\in\Fp^\times\quad\text{for any root }t_i\text{ of }f.
\end{equation}
Both sums are empty when $D_{\rm thin}=1$.
\end{corollary}
\begin{proof}
Subtracting the principal parts leaves a polynomial, which is zero because the
function and each summand vanish at infinity. Frobenius-conjugate poles have
conjugate residues; the residues lie in $\Fp$ by Theorem~\ref{thm:main}(e), so
they are equal on each irreducible factor. The identity
$f'/f=\sum_{f(t_i)=0}(T-t_i)^{-1}$ gives the second formula.
\end{proof}

\begin{theorem}[Frobenius on the eta quotient]\label{thm:C}
Let $\ell\ne p$ be prime. At every supersingular point $x$ of $X_0(\ell)_{\Fb}$,
\[
\Bigl(\frac{\Delta(\tau)}{\Delta(\ell\tau)}\Bigr)^{(p^2-1)/24}(x)=\leg{\ell}{p}.
\]
For $\ell=N\in\mathcal N$ this reads $u_N(x)^{(p^2-1)/a_N}=\Lp{N}$.
\end{theorem}

\begin{corollary}[a square root of Baker's congruence]\label{cor:D}
The value $c(E)$ is the same for all supersingular $E/\Fb$. Call it $c_p$. Then
$c_p\in\Fp$ and $c_p^2=-\Lp{-1}$; equivalently
\[
E_{p+1}^{(p-1)/2}\equiv c_p\,\Delta^{(p^2-1)/24}\pmod{(p,E_{p-1})}.
\]
\end{corollary}

\begin{conjecture}[numerical evaluation]\label{conj:E}
$c_p=\leg{6}{p}\cdot\bigl(\tfrac{p-1}{2}\bigr)!\bmod p$.
\end{conjecture}


\subsection{Consequences}\label{subsec:consequences}
The reconstruction in Corollary~\ref{cor:reconstruction} gives concrete formulas
once the supersingular values and their effective indices are known.
The connection to Nakaya's squarefree Fricke supersingular polynomials is a
change of coordinate: with the normalizations in \citet[\S4]{Nakaya19},
\begin{equation}\label{eq:nakaya-coordinate}
\mathcal S_{N,p}(T):=\prod_i(T-t_i)=ss_p^{(N*)}(T+c_N),\qquad
\begin{array}{c|rrrrr}
N&2&3&5&7&13\\\hline
c_N&104&42&16&9&0
\end{array}
\end{equation}
All constants are reduced modulo $p$; $\mathcal S_{N,p}$ is distinct from the
weighted polynomial $S$ of Proposition~\ref{prop:S}. Thus
$D_{\rm thin}=\mathcal S_{N,p}/\prod_{e_i>1}(T-t_i)$, while the omitted
values still occur in the weights $W_i$.
The exceptional values can also be specified by a polynomial independent of $p$.
Let $h_N$ be the elliptic-value polynomial in Table~\ref{tab:elliptic-values}.
Lemma~\ref{lem:elliptic-specialization} and Theorem~\ref{thm:main}(b) give
\begin{equation}\label{eq:elliptic-gcd}
D_{\rm thin}(T)=\frac{\mathcal S_{N,p}(T)}
{\gcd\bigl(\mathcal S_{N,p}(T),\overline h_N(T)\bigr)},
\end{equation}
where the gcd is monic. Thus the denominator is obtained from Nakaya's
supersingular polynomial by removing its common factors with the reduction
of a fixed characteristic-zero elliptic-value polynomial.

At level two, put
$\varepsilon_p=(1-\left(\frac{-1}{p}\right))/2$ and
$\nu_p=(1-\left(\frac{-2}{p}\right))/2$.
The level-two indices described in Section~\ref{sec:geometry} give
\begin{equation}\label{eq:leveltwo-denominator}
D_{2,p}(T)=\frac{ss_p^{(2*)}(T+104)}
{(T+104)^{\varepsilon_p}(T-152)^{\nu_p}}.
\end{equation}
The two removed values correspond to $0$ and $256$ in Nakaya's coordinate;
they are distinct for $p\ge5$.

\begin{corollary}[level-two degree and vanishing]\label{cor:leveltwo-vanishing}
For every prime $p\ge5$,
\begin{equation}\label{eq:leveltwo-degree}
\deg D_{2,p}=\frac{p-4+\left(\frac{-1}{p}\right)+2\left(\frac{-2}{p}\right)}8.
\end{equation}
In particular $T_{2A}|U_p\equiv0\pmod p$ if and only if $p=5$ or $7$.
\end{corollary}
\begin{proof}
The two possible thick values have indices four and two, and supersingularity
indicators $\varepsilon_p$ and $\nu_p$, respectively. Corollary~\ref{cor:degree}
therefore gives
$\deg D_{2,p}=(p-1)/8-\varepsilon_p/4-\nu_p/2$, which is
\eqref{eq:leveltwo-degree}. This degree is zero at $p=5,7$ and is at least
$(p-7)/8>0$ for $p\ge11$. Theorem~\ref{thm:main}(b) proves the last assertion.
\end{proof}

\citet[Lemma~4.8 and Table~4.3]{CMT19} establish $p$-adic
annihilation of $T_{2A}$ at $p=5,7$. The corollary specifies the exact
first-step modulo-$p$ vanishing criterion. Table~\ref{tab:leveltwo-examples}
illustrates the nonzero cases; at $p=17$, for example,
$R_{2,17}=D_{2,17}'/D_{2,17}$ and both residues are one.

\begin{table}[htbp]\centering
\begin{tabular}{rll}\toprule
$p$ & $D_{2,p}(T)$ & $P_{2,p}(T)$\\\midrule
11 & $T-5$ & $1$\\
13 & $T-3$ & $-2$\\
17 & $T^2-T-8$ & $2T-1$\\
29 & $T^3-12T^2+12T-6$ & $4T^2-5T-8$\\\bottomrule
\end{tabular}
\caption{Level-two examples, with $R_{2,p}=P_{2,p}/D_{2,p}$.
Coefficients use the unique integer representatives in $(-p/2,p/2)$.
These entries are included in the certificates of Section~\ref{sec:comp}.}
\label{tab:leveltwo-examples}
\end{table}

\subsection{Relation to earlier work}
The derivative argument in Section~\ref{sec:kz} follows the level-one method
presented by \citet[\S10]{KZ98}, who trace it to Koike's work
on Atkin's congruences \citep{Koike73}. The period formula of \citet{dS94}
and the congruence of \citet{Baker98} supply the level-one precedents for the
residue product and its sign; \citet{BL18} treat the corresponding Legendre
modular polynomials. The isogeny transformation of $B$ is due to
\citet[Th\'eor\`eme B, p.~106; \S3.3, Lemma~7 and (19), p.~124]{Robert80};
we include a $q$-expansion proof to fix the differential normalization.
Baker's congruence supplies the level-one factor $-\left(\frac{-1}{p}\right)$.

Fricke supersingular polynomials are studied by \citet{Sakai11,Sakai14},
\citet{Nakaya19} and \citet{Morton23}. These papers treat
orthogonal-polynomial representations, factorization and resultant
identities. Formula~\eqref{eq:nakaya-coordinate} identifies the supersingular
values used here with that literature; the present residue formula additionally
assigns a nonzero coefficient to each thin pole and retains the ramified values
in its weighted product.

\citet[\S4.2, Table~1 and Appendix~C]{JST19} determine the elliptic values
of normalized Fricke Hauptmoduln and their associated class fields. Their
$j_N$ equals our $T_N$. Table~\ref{tab:elliptic-values} records their
polynomials at the five levels considered here. Their Schwarzian equation
detects the same local ramification indices that enter our Hasse-derivative
calculation; \eqref{eq:elliptic-gcd} expresses its elliptic-value locus as
the locus of possible supersingular pole cancellations. The signed residue
formula concerns the remaining thin values.

\citet[Proposition~3.1 and \S4]{DS26} use Deligne's
supersingular expansion of $pJ|U_p$ and modular functions of prime-power level
to characterize the exponents in the order of the Monster. Their sharp
valuation bounds at the exceptional level-one invariants are closely related
to the ramification mechanism considered here. Our calculation determines the
signed residues of the reduced $U_p$ images on the five fixed Fricke curves.

Half-exponent discriminant identities have an earlier precedent in
\citet[Theorem~3, p.~72; \S4]{Landweber88}. For a supersingular curve
$y^2=4\prod_{a=1}^3(x-e_a)$ he proves
\[
\Delta^{(p^2-1)/24}=(e_a-e_b)^{(p^2-1)/4}\quad(a\ne b).
\]
Theorem~\ref{thm:C} instead compares discriminants along a prime-degree isogeny,
with its differential normalization built into a weight-zero function, and identifies
the resulting sign by an eta character. \citet[Theorem~1.2 and Remark~4]{Ono25}
proves the squared isogeny-ratio identity and records examples with nontrivial
half-exponent sign. His work develops the supersingular congruences arising in
the description of the divided beta family by \citet[Theorem~1.2]{Behrens09}
and the explicit calculations of \citet[\S5.4, Conjecture~5.17]{Larson20}.
These connections motivate the discriminant identity; no new statement about the
Adams--Novikov spectral sequence is needed here.

Our main result combines the signed isogeny-ratio identity with the Hasse-derivative
calculation to determine the full residue products on the five Fricke curves.
We make no claim that the exponent $(p^2-1)/24$ itself is new, or that the
isogeny transformation of $E_{p+1}$ originates here.

\section{Standard inputs and the residue criterion}\label{sec:imports}
We use the $q$-expansion principle and Kodaira--Spencer isomorphism for prime-to-$p$
level structures \citep{Katz73,KM85,Gross90}. On a fine moduli cover, the Hasse invariant
has simple zeros precisely at the supersingular points, and Kodaira--Spencer identifies
a weight-two form $f$ with $f\,dq/q$. In particular, $\theta T$ represents $dT$.
We record the statements used repeatedly.
\begin{fact}\label{fact:modp}
The form $A$ has $q$-expansion $1$ and simple supersingular zeros on the moduli stack.
The $q$-expansion principle holds at prime-to-$p$ level. For $p\ge5$ the automorphism
groups relevant to $X_0(N)$ and its Fricke quotient are tame. Coarse quotients
commute with reduction. The supersingular values of $T$ lie in $\Fq$.
\end{fact}
The last assertion follows by choosing a model of a supersingular curve with
$p^2$-Frobenius $[-p]$ (Lemma~\ref{lem:deuring}); every cyclic subgroup of order $N$
is then rational as a subgroup scheme.
\begin{fact}[Baker's congruence]\label{fact:baker}
For every $p\ge5$,
\[
E_{p+1}^{p-1}\equiv-\Lp{-1}\Delta^{(p^2-1)/12}\pmod{(p,E_{p-1})}.
\]
\end{fact}
This is \citet[Theorem~1.4]{Baker98}, with the normalized Eisenstein series used here.
\begin{fact}\label{fact:graph}
For every prime $\ell\ne p$, the supersingular $\ell$-isogeny graph over $\Fb$ is
connected \citep{Mestre86,Pizer90}.
\end{fact}

\begin{proposition}\label{prop:Psi}
Let $\Delta_p$ be the set of integral matrices of determinant $p$ with lower-left entry divisible by
$N$. There is a unique $\Psi_{N,p}\in\Z[X,Y]$ with
\[
\Psi_{N,p}(X,T(\tau))=(X-T(p\tau))\prod_{b=0}^{p-1}\Bigl(X-T\bigl(\tfrac{\tau+b}{p}\bigr)\Bigr).
\]
It is monic of degree $p+1$ in each variable, symmetric, and satisfies
$\Psi_{N,p}\equiv(X^p-Y)(X-Y^p)\pmod p$.
\end{proposition}
\begin{proof}
The matrices $\mathrm{diag}(p,1)$ and $\begin{pmatrix}1&b\\0&p\end{pmatrix}$, $0\le b<p$, represent
$\Gamma_0(N)\backslash\Delta_p$. Conjugation by $W_N=\begin{pmatrix}0&-1\\N&0\end{pmatrix}$ maps
$\begin{pmatrix}a&b\\c&d\end{pmatrix}$ to $\begin{pmatrix}d&-c/N\\-Nb&a\end{pmatrix}$, so it preserves
$\Delta_p$. Hence the elementary symmetric functions of the roots are $\Gamma_0(N)^+$-invariant,
holomorphic on the upper half-plane, with a pole only at the single cusp.
Their coefficients initially lie in $\Z[\zeta_p]((q^{1/p}))$.
Translation of $\tau$ removes fractional powers, and cyclotomic Galois automorphisms
permute the roots, so the expansions lie in $\Z((q))$.
Subtracting integral powers of $T=q^{-1}+O(q)$ then puts them in $\Z[T]$. Modulo $1-\zeta_p$ all $T((\tau+b)/p)$ coincide with $T(q^{1/p})$, whose $p$-th power is $T(q)$, and
$T(p\tau)\equiv T^p$; this gives the congruence. For irreducibility, the reduction of $\Gamma_0(N)$ modulo $p$ is
$\mathrm{SL}_2(\Fp)$, by the Chinese remainder theorem and the surjectivity of reduction
for $\mathrm{SL}_2(\Z)$. It acts transitively on the $p+1$ cyclic subgroups of order $p$.
The $p+1$ roots displayed above are distinct: their leading terms are $q^{-p}$ and
$\zeta_p^{-b}q^{-1/p}$ for $0\le b<p$. They therefore form one orbit, so
$\Psi_{N,p}(X,T)$ is irreducible over $\mathbb C(T)$.
The pole order of the $j$th elementary symmetric function is at most
$p+(j-1)/p\le p+1$; the constant coefficient has pole order $p+1$.
Thus the polynomial has degree $p+1$ in $Y$ as well as in $X$.
For symmetry, $\mathrm{diag}(1,p)\in\Delta_p$ gives
$\Psi_{N,p}(T(\tau),T(p\tau))=0$ by evaluating the defining relation at $p\tau$.
Irreducibility and Gauss's lemma therefore identify the reversed integral primitive
polynomial with the original up to sign.
Hence $\Psi_{N,p}(Y,X)=c\,\Psi_{N,p}(X,Y)$ with $c=\pm1$. Reducing modulo $p$, where both polynomials
equal the symmetric $(X^p-Y)(X-Y^p)$, forces $c=1$.
\end{proof}

Define $H=H_{N,p}\in\Z[X,Y]$ by $\Psi_{N,p}=(X^p-Y)(X-Y^p)+pH$; $H$ is symmetric. Put
\[
\psi=\frac{T(q)^p-T(q^p)}{p}\in\Z((q)),
\]
which is integral because $T(q^p)\equiv T(q)^p\pmod p$.

\begin{proposition}[residue criterion]\label{prop:res}
In $\Fp((q))$,
\[
\overline\psi=\frac{\overline H(T^p,T)}{T^{p^2}-T},\qquad
\overline{T|U_p}=\overline Q(T)+\overline\psi,
\]
where $\overline Q\in\Fp[T]$ is defined by $-[X^p]\Psi_{N,p}=Y^p+pQ(Y)$. Hence
$R_{N,p}=\overline Q(T)+\overline H(T^p,T)/(T^{p^2}-T)$ exists and is unique. It vanishes at
$\infty$, its poles are simple and lie in $\Fq$, and for $t\in\Fq$
\[
\Res_{T=t}R_{N,p}=-\overline H(t^p,t).
\]
\end{proposition}
\begin{proof}
Substitute $X=T(q^p)=T^p-p\psi$ and $Y=T$ into
$\Psi_{N,p}=0$. Since $X-Y^p=-p\psi$, division by $p$ and reduction give
$-\overline\psi\,(T^{p^2}-T)+\overline H(T^p,T)=0$. The sum of the roots of $\Psi_{N,p}(X,T)$ is
$T(q^p)+p\,(T|U_p)=T^p+pQ(T)$, so $T|U_p=Q(T)+\psi$. The polynomial $T^{p^2}-T$ is squarefree, its
roots are exactly $\Fq$, and its derivative is $-1$. $R_{N,p}$ vanishes at $\infty$ because
$T|U_p\in q\Z[[q]]$ (the constant term of $T$ is $0$). Uniqueness holds because $\overline T$ is
transcendental over $\Fp$.
\end{proof}

\begin{corollary}\label{cor:rFp}
For every $t\in\Fq$, $\Res_{t}R_{N,p}\in\Fp$ and $\Res_{t^p}R_{N,p}=\Res_tR_{N,p}$.
\end{corollary}
\begin{proof}
$(\Res_tR)^p=\Res_{t^p}R=-\overline H(t^{p^2},t^p)=-\overline H(t,t^p)=-\overline H(t^p,t)=\Res_tR$, using
$t^{p^2}=t$, the $\Fp$-rationality of $R$ and the symmetry of $H$.
\end{proof}

Proposition~\ref{prop:res} gives Theorem~\ref{thm:main}(a) and the a~priori bound
$\den R_{N,p}\mid T^{p^2}-T$ used by the certificates in Section~\ref{sec:comp}.

\section{Weights and the weighted supersingular polynomial}\label{sec:weights}

\begin{lemma}[effective inertia]\label{lem:inertia}
Let $x=(E,C)$ be a non-cuspidal point of $X_0(N)_{\Fb}$, supersingular or ordinary. Let $G_x$ be the group of automorphisms
of $(E,C)$ together with the isomorphisms $(E,C)\to w_N(E,C)=(E/C,E[N]/C)$ when $w_N(x)=x$. It is a group once
$w_N^2(E,C)\cong(E,C)$ is identified via $[N]$.
\begin{enumerate}[label=\textup{(\roman*)}]
\item $|G_x/\{\pm1\}|=m_x|\Stab_{\langle w_N\rangle}(x)|=e_x$. The value is the same at $x$ and at $w_N(x)$.
\item $\{\pm1\}$ is exactly the kernel of the action of $G_x$ on the deformation space of $x$. Since
$e_x\in\{1,2,3,4,6\}$ is prime to $p$, $G_x/\{\pm1\}$ is cyclic and acts faithfully on the tangent line.
Hence, for a local parameter $z$ of a fine moduli cover at a point over $x$, $T-t=z^{e_x}\cdot(\text{unit})$.
\item If $x$ is supersingular, $(p+1)/e_x\in\Z$.
\end{enumerate}
\end{lemma}
\begin{proof}
(i) If $w_N(x)=x$, the isomorphisms $(E,C)\to w_N(E,C)$ form a torsor under $\Aut(E,C)$, so
$|G_x|=2|\Aut(E,C)|$. Conjugation by the isogeny $E\to E/C$ identifies $\Aut(E,C)$ with
$\Aut(E/C,E[N]/C)$, so $m_x=m_{w_N(x)}$.
(ii) $[-1]$ acts trivially on deformations. Suppose $g\in G_x$, $g\ne\pm1$, acted trivially. Then $g$
would extend to the universal deformation. Its generic fibre has endomorphism ring $\Z$, and $g$
corresponds either to a nontrivial automorphism or to an endomorphism of degree $N$, which is not a square.
That is impossible. The same argument works at ordinary points. A finite group of order prime to $p$ acting faithfully on $k[[z]]$ is cyclic and
acts faithfully on $(z)/(z^2)$, so the invariant ring is $k[[z^{e}]]$.
(iii) $e=3$ or $6$ forces $m_x=3$, so $j=0$ is supersingular and $p\equiv2\pmod3$; $e=4$ forces
$m_x=2$, so $j=1728$ is supersingular and $p\equiv3\pmod4$. For $p\ge5$, $m_x\in\{1,2,3\}$.
\end{proof}

For explicit calculations, $m_x$ can be read off from the multiplicity $\mu$ of $u(x)$ as a root of
$P_N(u)-j\,u^N$: $m_x=3/\mu$ over $j=0$, $m_x=2/\mu$ over $j=1728$, and $m_x=1$ (with $\mu=1$) otherwise.
These are ramification indices of the coarse map $X_0(N)\to X(1)$.
The Fricke stabilizer has order two exactly when $u(x)^2=b_N$.

\begin{lemma}[elliptic values and specialization]\label{lem:elliptic-specialization}
Let $h_N(T)=\prod_{\tau}(T-T_N(\tau))$, where $\tau$ runs over the
inequivalent elliptic points of $\Gamma_0(N)^+$. For the five levels under
consideration, $h_N$ is given in Table~\ref{tab:elliptic-values}.
For every prime $p\ge5$ with $p\ne N$, its reduction is squarefree, and
its roots in $\Fb$ are exactly the values with effective index $e>1$.
Consequently,
\[
\gcd(\mathcal S_{N,p},\overline h_N)=\prod_{e_i>1}(T-t_i).
\]
\end{lemma}
\begin{proof}
The characteristic-zero values and elliptic orders are those of
\citet[\S4.2, Table~1 and Appendix~C]{JST19}, in the same normalization
$T_N=q^{-1}+O(q)$. Their elliptic orders are the effective indices of
Lemma~\ref{lem:inertia} in characteristic zero.
We justify the specialization of their support. A nontrivial effective
stabilizer arises either from an automorphism of $(E,C)$ or from a Fricke
fixed point. The former can occur only over $j=0$ or $1728$. Since
$p\ge5$ and $p\nmid N$, the corresponding elliptic curve and its
automorphisms lift to characteristic zero, and the prime-to-$p$ cyclic
subgroup $C$ lifts uniquely over a strict henselian base. An automorphism
preserving $C$ on the special fibre therefore preserves its lift.
Conversely, these prime-to-$p$ automorphisms persist under reduction.
Fricke fixed points are given by $u^2=b_N$, which is finite \'etale away
from $2N$ and has the same fixed-point locus under specialization.
This accounts for all values with $e>1$ in either characteristic.
Directly from the displayed polynomials, the discriminants of $h_N$, in
the order $N=2,3,5,7,13$, are
\[
2^{16},\qquad 2^4\cdot3^6,\qquad 2^{12}\cdot5^3,\qquad
2^4\cdot3^6\cdot7^2,\qquad 2^{12}\cdot3^6\cdot13.
\]
Thus distinct elliptic values do not collide at any admissible prime.
The roots of $\overline h_N$ are precisely the values with $e>1$;
intersecting this set with the supersingular locus proves the gcd formula.
\end{proof}

\begin{table}[htbp]\centering
\begin{tabular}{cll}\toprule
$N$ & $h_N(T)$ & Effective indices\\\midrule
2 & $(T+104)(T-152)$ & $4,\ 2$\\
3 & $(T+42)(T-66)$ & $6,\ 2$\\
5 & $(T+16)(T^2-12T-464)$ & $2,\ 2$\\
7 & $(T+10)(T+9)(T-18)$ & $2,\ 3,\ 2$\\
13 & $(T+4)(T+3)(T^2-4T-48)$ & $2,\ 3,\ 2$\\\bottomrule
\end{tabular}
\caption{Elliptic-value polynomials from \citet{JST19}.
The indices are listed in the order of the displayed factors; an index
beside a quadratic factor applies to both roots. At admissible primes,
the supersingular roots of $\overline h_N$ are exactly the thick values
removed in \eqref{eq:elliptic-gcd}.}
\label{tab:elliptic-values}
\end{table}

\begin{lemma}\label{lem:AW}
As mod-$p$ forms on $\Gamma_0(N)$: $A|_{p-1}W_N=\Lp{N}A$, $\Delta_N|_{24}W_N=\Delta_N$ and
$\theta T|_2W_N=\theta T$.
\end{lemma}
\begin{proof}
For a level-one form $f$ of weight $k$, $f|_kW_N=N^{k/2}f(N\tau)$. The form $E_{p-1}(N\tau)$ has
$q$-expansion $\equiv1\equiv E_{p-1}(\tau)$, so it equals $A$ modulo $p$ by Fact~\ref{fact:modp}(i);
and $N^{(p-1)/2}\equiv\Lp N$. The other two identities hold in characteristic $0$:
$\Delta(-1/(N\tau))\Delta(-1/\tau)=N^{12}\tau^{24}\Delta(N\tau)\Delta(\tau)$, and $T$ is
$w_N$-invariant.
\end{proof}

\begin{proposition}[weighted supersingular polynomial]\label{prop:S}
Let $S(T)=\prod_{k=1}^s(T-t_k)^{(p+1)/e_k}$. Then $S\in\Fp[T]$ and
\[
S(\overline T)=\Delta_N^{-k}=q^{-(N+1)k}\prod_{n\ge1}(1-q^n)^{-24k}(1-q^{Nn})^{-24k}\quad\text{in }\Fp((q)).
\]
Equivalently, $S=A^{p+1}/\Delta_N^{k}$ as functions on $X_0(N)^+_{\Fp}$. In particular
$\sum_k e_k^{-1}=(N+1)(p-1)/24$.
\end{proposition}
\begin{proof}
$F:=A^{p+1}/\Delta_N^k$ is a ratio of mod-$p$ forms of weight $p^2-1$. By Lemma~\ref{lem:AW} it is
$w_N$-invariant, since $\Lp N^{p+1}=1$. $\Delta_N$ does not vanish on $Y_0(N)$, so $F$ is regular off the cusp of
$X_0(N)^+$. Its $q$-expansion there is $\Delta_N^{-k}=q^{-(N+1)k}(1+O(q))$, because $A$ has
$q$-expansion $1$. Hence $F$ is a monic polynomial in $T$ of degree $(N+1)k$. By
Fact~\ref{fact:modp}(ii) and Lemma~\ref{lem:inertia}(ii), $F$ vanishes exactly at the supersingular
$t_k$, to order $(p+1)/e_k$ in $T-t_k$. Comparing degrees gives the mass identity.
\end{proof}

\begin{corollary}\label{cor:degree}
The degree of the reduced denominator is
\[
\deg D_{\rm thin}=\frac{(N+1)(p-1)}{24}-\sum_{e_i>1}\frac1{e_i}.
\]
In particular $R_{N,p}=0$ exactly when no thin supersingular value exists.
\end{corollary}
\begin{proof}
Separate the thin and non-thin terms in the mass identity of
Proposition~\ref{prop:S}; use Theorem~\ref{thm:main}(b) for the last assertion.
\end{proof}

\section{The Kaneko--Zagier step}\label{sec:kz}

The following is the level-$N$ version of \citet[\S10, (35)--(38)]{KZ98}, who credit the method to Koike.

\begin{lemma}\label{lem:dpsi}
In $\Fp((q))$, with $d/dT=(\theta T)^{-1}\theta$,
\[
\frac{d\overline\psi}{dT}=T^{p-1}-(\theta T)^{p-1},\qquad (\theta T)^{p-1}=g(\overline T),\quad
g:=\frac{(\theta T)^{p-1}}{A^2}\in\Fp(T).
\]
\end{lemma}
\begin{proof}
Over $\Z$, $\theta\psi=T^{p-1}\theta T-(\theta T)(q^p)$, and $(\theta T)(q^p)\equiv(\theta T)^p$.
Divide by $\theta T=-q^{-1}+\cdots$, which is invertible. The function $g$ is a ratio of mod-$p$
forms of weight $2(p-1)$, it is $w_N$-invariant by Lemma~\ref{lem:AW}, and it has $q$-expansion
$(\theta T)^{p-1}$ because $A\equiv1$.
\end{proof}

\begin{lemma}[local structure of $g$]\label{lem:g}
Let $t\in\Fq$ be a non-cuspidal point with effective inertia $e_t$, and put $\varepsilon_t=1$ if
$t$ is supersingular and $0$ otherwise. Then
\[
\mathrm{ord}_{T=t}\,g=\frac{(p-1)(e_t-1)-2\varepsilon_t}{e_t}.
\]
For thin $t$ and $x$ over $t$, let $a=A/(\theta T)^{(p-1)/2}$, a function on $X_0(N)$ near $x$.
Then $a$ has a simple zero at $x$,
\[
a'(t):=\frac{da}{dT}(x)=\frac{(B/12)(x)}{(\theta T)(x)^{(p+1)/2}}\ne0,
\qquad g=a'(t)^{-2}(T-t)^{-2}\bigl(1+O(T-t)\bigr).
\]
\end{lemma}
\begin{proof}
By Fact~\ref{fact:modp}(iv) and Lemma~\ref{lem:inertia}(ii), $\theta T$ vanishes to order $e_t-1$ in a
stack parameter $z$, and $A$ to order $\varepsilon_t$; also $T-t\sim z^{e_t}$. This gives the order formula.
For thin $t$, the map $X_0(N)\to X_0(N)^+$ is \'etale at $x$ and $\theta T(x)\ne0$. In
characteristic $0$, for forms $F$ and $M$ of equal weight,
$\theta(F/M)=(\partial F\cdot M-F\cdot\partial M)/M^2$, where $\partial=\theta-\frac{k}{12}E_2$ is the
Serre derivative. Apply this with $F=E_{p-1}$ and $M=(\theta T)^{(p-1)/2}$ and reduce modulo $p$; at
$x$ we have $A(x)=0$. The form $\partial E_{p-1}$ has weight $p+1$ and $q$-expansion
$\equiv\theta(1)+\frac1{12}E_2\equiv\frac1{12}E_{p+1}$, using the Kummer congruence
$2(p+1)/B_{p+1}\equiv24$ and $\sigma_p\equiv\sigma_1$. So $\partial A=B/12$ as mod-$p$ forms. Finally
$da/dT=\theta a/\theta T$. The value $B(x)\ne0$ because $B/12$ is the derivative of $A$ at its simple zero.
\end{proof}

\begin{proof}[Proof of Theorem~\ref{thm:main}(b)--(c)]
Let $t\in\Fq$. By Proposition~\ref{prop:res}, near $t$ we have
$\overline\psi=-\overline H(t^p,t)(T-t)^{-1}+O(1)$, because $T^{p^2}-T=-(T-t)(1+O(T-t))$. So the
coefficient of $(T-t)^{-2}$ in $d\overline\psi/dT$ is $\overline H(t^p,t)$. By Lemma~\ref{lem:dpsi}
it is also $-[g]_{-2}$. Hence
\[
\Res_tR_{N,p}=-\overline H(t^p,t)=[g]_{-2}.
\]
By Lemma~\ref{lem:g}, $\mathrm{ord}_tg\ge0$ unless $t$ is thin. For $p\ge5$ and $e_t\ge2$ we have
$(p-1)(e_t-1)-2\ge0$. At thin $t$, $[g]_{-2}=a'(t)^{-2}=144(\theta T)(x)^{p+1}/B(x)^2\ne0$. Every pole is
simple by Proposition~\ref{prop:res}, so $R_{N,p}=P/D_{\rm thin}$ with $P(t_i)\ne0$ at every thin
$t_i$. Since $R_{N,p}(\infty)=0$, $\deg P<\deg D_{\rm thin}$.
\end{proof}

\begin{lemma}[Robert's isogeny identity]\label{lem:Bisog}
Let $N$ be any prime $\ne p$ (not only $N\in\mathcal N$). Let $x=(E,C)$ be a supersingular point of
$X_0(N)$, $\pi:E\to E'=E/C$, and $\omega=\pi^*\omega'$. Then
$B(E',\omega')=N\,B(E,\omega)$, $A(E',\omega')=A(E,\omega)$ and
$\Delta(N\tau)(x)=N^{-12}\Delta(E',\omega')$.
\end{lemma}
\begin{proof}
The isogeny identity is \citet[Th\'eor\`eme B; \S3.3, Lemma~7]{Robert80}; see also
\citet[Lemma~2.2]{Baker98}. We verify the normalization by $q$-expansions.
For a level-one form $f$ of weight $w$, $f(N\tau)$ at $(E,C,\omega)$ equals $N^{-w}f(E',\omega')$. Check
this on $(\mathrm{Tate}(q),\mu_N)$, where $E'\cong\mathrm{Tate}(q^N)$ and
$\pi^*\omega_{\rm can}'=N\omega_{\rm can}$. The form $E_2(\tau)-NE_2(N\tau)$ is a holomorphic
weight-$2$ form on $\Gamma_0(N)$ with $p$-integral coefficients. Modulo $p$, the forms
$B(\tau)-NB(N\tau)$ and $A\cdot(E_2-NE_2(N\tau))$ have weight $p+1$ and the same $q$-expansion, so they
are equal. At a supersingular point this gives $B(E,\omega)=N\cdot N^{-(p+1)}B(E',\omega')=N^{-1}B(E',\omega')$,
because $N^{p-1}=1$. The statement about $A$ follows in the same way from $A(N\tau)=A(\tau)$.
\end{proof}

\begin{proposition}[weighted residue product]\label{prop:product}
For thin $t_i$ and $x=(E,C)$ over $t_i$,
\[
r_iW_i=\frac{B^{p-1}}{\Delta_N^{\,k}}(x)=\beta(E)\Bigl(\frac{\Delta(\tau)}{\Delta(N\tau)}\Bigr)^{k}(x)
=\Lp{N}c(E)c(E'),\qquad \beta(E):=\frac{B(E,\omega)^{p-1}}{\Delta(E,\omega)^{2k}}=c(E)^2.
\]
\end{proposition}
\begin{proof}
Near $x$, Proposition~\ref{prop:S} gives $S(T)=A^{p+1}/\Delta_N^k=a^{p+1}\,Q^{k}$ with
$Q:=(\theta T)^{12}/\Delta_N$, a function that is regular and nonzero at $x$. Hence
$W_i=[S/(T-t_i)^{p+1}]_{T=t_i}=a'(t_i)^{p+1}Q(x)^k$. With $r_i=a'(t_i)^{-2}$ this gives
\[
r_iW_i=a'(t_i)^{p-1}Q(x)^k=\Bigl(\frac{B}{12}\Bigr)^{p-1}(\theta T)^{-\frac{p^2-1}2}\cdot
\frac{(\theta T)^{\frac{p^2-1}{2}}}{\Delta_N^{k}}=\frac{B^{p-1}}{\Delta_N^k}(x).
\]
Write $\Delta_N^k=\Delta^{2k}(\Delta(N\tau)/\Delta)^k$ for the second equality. For the third,
Lemma~\ref{lem:Bisog} gives $\Delta_N(x)=N^{-12}\Delta(E,\omega)\Delta(E',\omega')$ with
$N^{12k}=N^{(p^2-1)/2}=1$, and $B(E)^{(p-1)/2}=N^{-(p-1)/2}B(E')^{(p-1)/2}=\Lp NB(E')^{(p-1)/2}$.
\end{proof}

\section{The eta quotient at supersingular points}\label{sec:eta}

Throughout this section $\ell\ne p$ is a prime, $M:=576\ell$, $k:=(p^2-1)/24$, and
\[
D_\ell(\tau):=\frac{\Delta(\tau)}{\Delta(\ell\tau)},\qquad
G(\tau):=\frac{\eta(24\ell\tau)}{\eta(24\tau)}.
\]
Then $G=q^{\ell-1}\prod_{n\ge1}(1-q^{24\ell n})/(1-q^{24n})\in q^{\ell-1}\Z[[q]]$
and $G^{24}=D_\ell(24\tau)^{-1}$.
We prove Theorem~\ref{thm:C}: $D_\ell(x)^{k}=\leg{\ell}{p}$ at every supersingular point $x$ of
$X_0(\ell)_{\Fb}$. The argument has four steps: (i) the multiplier of $G$ on $\Gamma_0(M)$; (ii) the
moduli interpretation of $D_\ell$, of $G$, and of the diamond operators; (iii) the Frobenius action
on a lift of $x$ to $X_1(M)$; (iv) the descent.

\subsection{The multiplier of \texorpdfstring{$G$}{G}}

We use the closed form of the eta multiplier (Rademacher; Petersson; see
\citealp[Ch.~4, \S4.1]{Knopp70}, and \citealp[Thm.~3.4]{Apostol90} for the Dedekind-sum form):
for $\gamma=\begin{pmatrix}a&b\\c&d\end{pmatrix}\in\mathrm{SL}_2(\Z)$ with $c>0$,
\begingroup\small
\begin{equation}\label{eq:petersson}
\eta(\gamma\tau)=\nu(\gamma)\,(c\tau+d)^{1/2}\,\eta(\tau),\qquad
\nu(\gamma)=
\begin{cases}
\leg{d}{c}\,e\!\Bigl(\dfrac{(a+d)c-bd(c^2-1)-3c}{24}\Bigr), & c\ \text{odd},\\[2ex]
\leg{c}{d}\,e\!\Bigl(\dfrac{(a+d)c-bd(c^2-1)+3d-3-3cd}{24}\Bigr), & c\ \text{even},
\end{cases}
\end{equation}
\endgroup
where $e(x)=e^{2\pi ix}$, $(c\tau+d)^{1/2}$ is the principal branch, and the symbols are Kronecker
symbols (so $\leg{c}{d}$ is defined for $d<0$, with $\leg{c}{-1}=\operatorname{sgn}c$). 

\begin{lemma}[multiplier of $G$]\label{lem:mult}
Let $\chi_\ell(d):=\leg{\ell}{d}$ (Kronecker symbol). For every
$\gamma=\begin{pmatrix}a&b\\c&d\end{pmatrix}\in\Gamma_0(M)$,
\[
G(\gamma\tau)=\chi_\ell(d)\,G(\tau).
\]
The function $d\mapsto\chi_\ell(d)$ on $(\Z/M)^\times$ is the Kronecker character of
$\Q(\sqrt\ell)$: it is a Dirichlet character of conductor $\ell$ ($\ell\equiv1\bmod4$),
$4\ell$ ($\ell\equiv3\bmod4$) or $8$ ($\ell=2$), each dividing $M$; it is even and quadratic, and
$\chi_\ell(p)=\leg{\ell}{p}$ is the Legendre symbol.
\end{lemma}
\begin{proof}
If $c=0$ then $\gamma=\pm\begin{pmatrix}1&b\\0&1\end{pmatrix}$, $G(\tau+b)=G(\tau)$ because $G$ has
an integral $q$-expansion, and $\chi_\ell(\pm1)=1$. If $c<0$ replace $\gamma$ by $-\gamma$, which
acts in the same way on $\mathfrak H$ and has $\chi_\ell(-d)=\chi_\ell(d)$ since $\leg{\ell}{-1}=1$.
So let $c>0$. Put $c_1=c/24$, $c_2=c/(24\ell)$, $b_1=24b$, $b_2=24\ell b$ and
$\gamma_i=\begin{pmatrix}a&b_i\\c_i&d\end{pmatrix}\in\mathrm{SL}_2(\Z)$. Then
$24\gamma\tau=\gamma_1(24\tau)$ and $24\ell\gamma\tau=\gamma_2(24\ell\tau)$, and both automorphy
factors equal $(c\tau+d)^{1/2}$, so $G(\gamma\tau)=\bigl(\nu(\gamma_2)/\nu(\gamma_1)\bigr)G(\tau)$.
Since $576\ell\mid c$, both $c_1$ and $c_2$ are divisible by $24$, in particular even, so the second
line of \eqref{eq:petersson} applies to both. In the quotient of the exponentials every term is a
multiple of $24$: $(a+d)(c_2-c_1)$ and $3d(c_2-c_1)$ because $24\mid c_i$, and
$d\bigl(b_2(c_2^2-1)-b_1(c_1^2-1)\bigr)$ because $24\mid b_i$. Hence the exponentials cancel and
\[
\frac{\nu(\gamma_2)}{\nu(\gamma_1)}=\leg{c_2}{d}\leg{c_1}{d}=\leg{c_1c_2}{d}=\leg{\ell c_2^2}{d}
=\leg{\ell}{d}\leg{c_2}{d}^{2}=\leg{\ell}{d},
\]
using $c_1=\ell c_2$ and $\gcd(c_2,d)=1$. The statements about $\chi_\ell$ are quadratic
reciprocity: for $d>0$ odd, $\leg{\ell}{d}=\leg{d}{\ell}$ if $\ell\equiv1\pmod 4$,
$\leg{\ell}{d}=\leg{-1}{d}\leg{d}{\ell}$ if $\ell\equiv3\pmod4$, and $\leg{2}{d}=\chi_8(d)$; in each
case this is the Kronecker character of discriminant $\ell$, $4\ell$, $8$ respectively, which is even
because $\Q(\sqrt\ell)$ is real. For a prime $p\nmid M$, $\chi_\ell(p)$ is the Legendre symbol.
\end{proof}

\begin{remark}
The multiplier lemma is also an instance of the Gordon--Hughes--Newman
eta-quotient transformation law \citep[Theorem~1.64]{Ono04}.
For $y=\eta(\tau)/\eta(\ell\tau)$, the corresponding formula requires
$24\ell\mid c$ and $24\mid b$; the latter hypothesis cannot in general be omitted.
\end{remark}

\subsection{Moduli conventions}\label{subsec:moduli}

\paragraph{The $\mu_M$-model.} For $M\ge5$ and over $\Z[1/M]$ let $Y_1(M)$ be the fine, affine moduli
scheme of pairs $(E,\iota)$ with $\iota:\mu_M\hookrightarrow E$ a closed immersion of group schemes,
and $X_1(M)\supset Y_1(M)$ its smooth proper compactification
\citep[IV--V]{DR73}; \citep[\S2]{Gross90}; \citep[\S1]{Katz76}. The Tate curve
$\mathrm{Tate}(q)=\mathbb G_m/q^{\Z}$ over $\Z[1/M]((q))$ with the canonical embedding
$\mu_M\subset\mathbb G_m$ gives a section $\infty$ of $X_1(M)$ over $\Z[1/M]$ whose formal
completion is $\operatorname{Spf}\Z[1/M][[q]]$. The complex points are $\Gamma_1(M)\backslash\mathfrak H^*$
via $\tau\mapsto x_\tau:=(E_\tau,\iota_\tau)$, $E_\tau=\mathbb C/(\Z+\Z\tau)$,
$\iota_\tau(e^{2\pi i/M})=1/M$; under $E_\tau\cong\mathbb C^\times/q^\Z$, $z\mapsto e^{2\pi iz}$,
this is the canonical embedding, so the $q$-expansion of a function at $\infty$ is its Fourier
expansion in $q=e^{2\pi i\tau}$. The diamond operators are the $\Z[1/M]$-automorphisms
\[
\langle d\rangle(E,\iota)=(E,\iota\circ[d]),\qquad d\in(\Z/M)^\times .
\]
\emph{Convention.} We write $\langle d\rangle$ for this automorphism of the scheme and
$f\circ\langle d\rangle$ for the pull-back of a function; no inverse occurs anywhere below.

\begin{lemma}[analytic diamond operators]\label{lem:diamond}
For $\gamma=\begin{pmatrix}a&b\\c&d\end{pmatrix}\in\Gamma_0(M)$ one has
$x_{\gamma\tau}=\langle d\rangle x_\tau$. Consequently, for a function $f$ on $X_1(M)_{\mathbb C}$
with $f(\gamma\tau)=\chi(d)f(\tau)$ on $\Gamma_0(M)$,
\[
f\circ\langle d\rangle=\chi(d)\,f .
\]
\end{lemma}
\begin{proof}
Multiplication by $c\tau+d$ maps $\Z+\Z\gamma\tau$ onto $\Z+\Z\tau$ (it sends $\gamma\tau\mapsto a\tau+b$
and $1\mapsto c\tau+d$, and $\det\gamma=1$), hence is an isomorphism $E_{\gamma\tau}\to E_\tau$. It
sends $\iota_{\gamma\tau}(\zeta_M)=1/M$ to $(c\tau+d)/M\equiv d/M=\iota_\tau(\zeta_M^d)$ because
$M\mid c$. So $(E_{\gamma\tau},\iota_{\gamma\tau})\cong(E_\tau,\iota_\tau\circ[d])$. Then
$(f\circ\langle d\rangle)(x_\tau)=f(x_{\gamma\tau})=f(\gamma\tau)=\chi(d)f(x_\tau)$.
\end{proof}

\begin{lemma}[$q$-expansion descent]\label{lem:qexp}
Let $f\in\mathbb C(X_1(M))$ have poles supported at the cusps and Fourier expansion
at $\infty$ in $\Q((q))$. Then $f\in\Q(X_1(M))$ for the $\mu_M$-model.
\end{lemma}
\begin{proof}
The reduced cusp divisor is defined over $\Q$. Choose a multiple $D$ of it bounding
the poles of $f$. The space $H^0(X_1(M),\mathcal O(D))$ is defined over $\Q$, and its
base change to $\mathbb C$ contains $f$. Expansion at the rational cusp, in the
rational parameter $q$, gives a linear map defined over $\Q$. Finitely many Laurent
coefficients give an injective map on this finite-dimensional space: otherwise a
nonzero section could vanish to arbitrarily high order at $\infty$.
In a rational basis this injective coefficient map is a rational matrix. The
specified coefficients of $f$ are rational, so its coordinates in that basis are
rational. This proves descent from $\mathbb C$, including algebraicity.
\end{proof}

\begin{remark}[the point convention]\label{rem:point}
For the model classifying $(E,P)$ with $P$ of order $M$, the standard Tate cusp
$(\mathrm{Tate}(q),\zeta_M)$ is defined over $\Q(\zeta_M)$ rather than $\Q$
\citep[Example~2.3]{ConradQexp}. Rational Fourier coefficients do not by themselves
give the same descent assertion on that model. We use only the $\mu_M$-model,
its rational Tate cusp and the diamond convention fixed above.
\end{remark}

\paragraph{The function $D_\ell$ on $Y_0(\ell)$.} For a pair $(E,C)$ over a $\Z[1/\ell]$-scheme, with
$\pi:E\to E'=E/C$ the quotient isogeny and $\omega'$ any nowhere-vanishing invariant differential on
$E'$, put
\begin{equation}\label{eq:Dmoduli}
D_\ell(E,C):=\ell^{12}\,\frac{\Delta(E,\pi^*\omega')}{\Delta(E',\omega')}\in\mathcal O(S)^\times .
\end{equation}
This is independent of $\omega'$ (both discriminants have weight $12$), so it is an invertible
function on $Y_0(\ell)_{\Z[1/\ell]}$. On the Tate curve with $C=\mu_\ell$, $E'=\mathrm{Tate}(q^\ell)$
via $t\mapsto t^\ell$, $\pi^*(dt'/t')=\ell\,dt/t$, so
$\Delta(E,\pi^*\omega'_{\rm can})=\ell^{-12}\Delta(q)$ and $\Delta(E',\omega'_{\rm can})=\Delta(q^\ell)$:
the $q$-expansion of \eqref{eq:Dmoduli} is $\Delta(q)/\Delta(q^\ell)$. Hence
\eqref{eq:Dmoduli} is the algebraic incarnation of $\Delta(\tau)/\Delta(\ell\tau)$, with the value
at a geometric point $x=(E,C)$ given by \eqref{eq:Dmoduli}; this is the meaning of
``$D_\ell(x)$'' in Theorem~\ref{thm:C} (with the differential normalization of Lemma~\ref{lem:Bisog}).

\begin{lemma}[the degeneracy map]\label{lem:delta}
The rule
\[
\delta(E,\iota)=\Bigl(E/\iota(\mu_{24}),\ \ \iota(\mu_{24\ell})/\iota(\mu_{24})\Bigr)
\]
defines a morphism $\delta:Y_1(M)\to Y_0(\ell)$ over $\Z[1/M]$ (here $\mu_{24}\subset\mu_{24\ell}\subset\mu_M$).
On complex points $\delta(x_\tau)=(E_{24\tau},\langle 1/\ell\rangle)$, i.e.\ $\delta$ is
$\tau\mapsto24\tau$, and
\[
G^{24}=\delta^*\bigl(D_\ell^{-1}\bigr)\quad\text{in }\Q(X_1(M)).
\]
For every $x\in Y_0(\ell)(\Fb)$ there is $\tilde x\in Y_1(M)(\Fb)$ with $\delta(\tilde x)=x$.
\end{lemma}
\begin{proof}
The quotient by the finite \'etale subgroup $\iota(\mu_{24})$ and the image of $\iota(\mu_{24\ell})$
(cyclic of order $\ell$) are defined over $\Z[1/M]$. Analytically, multiplication by $24$ induces
$E_\tau\to\mathbb C/(\Z+\Z\,24\tau)=E_{24\tau}$ with kernel $\langle 1/24\rangle=\iota_\tau(\mu_{24})$,
and maps $\iota_\tau(\mu_{24\ell})=\langle1/(24\ell)\rangle$ onto $\langle1/\ell\rangle$. The
function $D_\ell$ at $(E_{\tau'},\langle1/\ell\rangle)$ is $\Delta(\tau')/\Delta(\ell\tau')$, so
$(\delta^*D_\ell)(x_\tau)=\Delta(24\tau)/\Delta(24\ell\tau)=G(\tau)^{-24}$; Lemma~\ref{lem:qexp} applies to $G$, whose poles are cuspidal and whose Fourier
coefficients are integral; thus $G\in\Q(X_1(M))$. Equality on complex points now
gives the asserted identity in that rational function field. For surjectivity on $\Fb$-points, choose a cyclic subgroup $K\subset E$ of
order $24$, requiring $K[\ell]\ne C$ when $\ell=2$ or $3$.
There is such a choice because the $\ell$-torsion has $\ell+1$ lines, and each line
extends to a cyclic subgroup of the required prime-power order.
Let $q_K:E\to E_1=E/K$ and let $\widehat q_K:E_1\to E$ be its dual, with
kernel $K_1$. Elementary divisors of prime-to-$p$ torsion show that
$H=\widehat q_K^{-1}(C)$ is cyclic of order $24\ell$.
For clarity, at a prime $r\mid24$, choose torsion coordinates in which $K$
is the first coordinate line. The dual quotient has elementary divisors
$(1,r^{v_r(24)})$. If $r=\ell$, the condition $C\ne K[r]$ makes its preimage
cyclic of order $r^{v_r(24)+1}$; at $r\ne\ell$ it is the cyclic kernel.
When $\ell\nmid24$, the additional $\ell$-part is a line and the same conclusion follows.
Extend $H$ to a cyclic subgroup $Z\subset E_1$ of order $M=24^2\ell$.
Its unique subgroup of order $24$ is $K_1$. Choose an isomorphism
$\iota:\mu_M\xrightarrow{\sim} Z$ over $\Fb$.
The dual quotient identifies $E_1/K_1$ with $E$, and
$\iota(\mu_{24\ell})/K_1$ with $C$. Thus $\delta(E_1,\iota)=x$.
\end{proof}

\begin{proposition}[descent and integrality of $G$]\label{prop:Gint}
$G\in\Q(X_1(M))$, and $G$ is an invertible regular function on $Y_1(M)_{\Z_{(p)}}$. Its reduction
$\overline G$ is an invertible regular function on $Y_1(M)_{\Fp}$, defined over $\Fp$, and
\[
\overline G\circ\langle d\rangle=\chi_\ell(d)\,\overline G\qquad(d\in(\Z/M)^\times),\qquad
\overline G^{\,24}=\delta^*\bigl(\overline D_\ell^{\,-1}\bigr).
\]
\end{proposition}
\begin{proof}
$G$ is holomorphic and non-vanishing on $\mathfrak H$, invariant under $\Gamma_1(M)$
(Lemma~\ref{lem:mult}, $\chi_\ell(1)=1$) and meromorphic at the cusps, so $G\in\mathbb C(X_1(M))$;
by Lemma~\ref{lem:qexp} and the integral $q$-expansion, $G\in\Q(X_1(M))$. By
Lemma~\ref{lem:delta}, $F:=\delta^*(D_\ell^{-1})$ lies in $\mathcal O(Y_1(M)_{\Z[1/M]})^\times$ and
$G^{24}=F$. The ring $R:=\mathcal O(Y_1(M)_{\Z_{(p)}})$ is a normal domain (smooth over $\Z_{(p)}$,
flat with integral generic fibre), and $G\in\operatorname{Frac}(R)$ satisfies the monic equation
$T^{24}-F=0$ with $F\in R^\times$; hence $G\in R$, and $G^{-1}=G^{23}F^{-1}\in R$. The identity
$G\circ\langle d\rangle=\chi_\ell(d)G$ holds in $\mathbb C(X_1(M))$ by Lemmas~\ref{lem:mult}
and~\ref{lem:diamond}, hence in $R$ ($\langle d\rangle$ is an automorphism of $Y_1(M)$ over
$\Z[1/M]$), hence after reduction. Reduction of $G^{24}=F$ gives the last identity.
\end{proof}

\subsection{Frobenius on a lift}

\begin{lemma}[Deuring models]\label{lem:deuring}
Let $E_1/\Fb$ be supersingular, $p\ge5$. Then $E_1$ has a model over $\Fq$ whose $p^2$-power
Frobenius endomorphism is $[-p]$. For such a model every finite subgroup of order prime to $p$ is
$\Fq$-rational, and $E_1/K$ again has Frobenius $[-p]$.
\end{lemma}
\begin{proof}
For any model $E/\Fq$ of a supersingular curve, the $p$-power Frobenius $F:E\to E^{(p)}$ and the
Verschiebung $V:E^{(p)}\to E$ are both purely inseparable of degree $p$ (the Hasse invariant
vanishes), so $F^{(p)}:E^{(p)}\to E^{(p^2)}=E$ and $V$ differ by an automorphism $u$ of $E$:
$F^{(p)}=u\circ V$, whence $\pi_{p^2}=F^{(p)}\circ F=u\circ V\circ F=u\circ[p]$. Thus
$\pi_{p^2}=p\,u$ with $u\in\Aut(E)$. If $j\ne0,1728$ then $u=\pm1$ and the quadratic twist changes
the sign; if $j=0$ (resp.\ $1728$) the curve $y^2=x^3+1$ (resp.\ $y^2=x^3+x$) over $\Fp$ has
trace $0$, so $\pi_{p^2}=\pi_p^2=[-p]$ \citep{Deuring41,Waterhouse69}; see also
\citet[Lemma~4.3]{Ono25}. Since $[-p]$ is central, Galois preserves every subgroup, and the quotient
isogeny $\phi$ satisfies $\pi'\phi=\phi\pi=\phi[-p]=[-p]\phi$, so $\pi'=[-p]$.
\end{proof}

\begin{proof}[Proof of Theorem~\ref{thm:C}]
Let $x=(E,C)\in Y_0(\ell)(\Fb)$ be supersingular and choose, by Lemma~\ref{lem:delta}, a lift
$\tilde x=(E_1,\iota)\in Y_1(M)(\Fb)$ with $\delta(\tilde x)=x$. The curve $E_1$ is isogenous to $E$,
hence supersingular, so by Lemma~\ref{lem:deuring} we may replace $(E_1,\iota)$ by an isomorphic
pair (same geometric point $\tilde x$, same value $\overline G(\tilde x)$) in which $E_1$ is defined
over $\Fq$ with $\pi_{E_1}=[-p]$. The embedding $\iota:\mu_M\to E_1$ is defined over $\Fb$; nothing
requires it to be $\Fq$-rational.

\emph{Galois action on the point.} Let $\sigma\in\operatorname{Gal}(\Fb/\Fq)$ be $u\mapsto u^{p^2}$.
Since $Y_1(M)$ is a fine moduli scheme, $\sigma$ acts on $Y_1(M)(\Fb)$ by
$\sigma(E_1,\iota)=(E_1^\sigma,\iota^\sigma)=(E_1,\sigma_{E_1}\circ\iota\circ\sigma_{\mu_M}^{-1})$,
where $\sigma_{\mu_M}(\zeta)=\zeta^{p^2}$ and $\sigma_{E_1}$ is the action on $E_1(\Fb)$ through the
$\Fq$-structure, i.e.\ the Frobenius endomorphism $\pi_{E_1}=[-p]$. Hence
\[
\iota^\sigma(\zeta)=[-p]\,\iota(\zeta^{p^{-2}})=[-p][p^{-2}]\,\iota(\zeta)=\iota([-p^{-1}]\zeta),
\qquad\text{i.e.}\qquad \sigma(\tilde x)=\langle -p^{-1}\rangle\tilde x,
\]
with $-p^{-1}$ read in $(\Z/M)^\times$. (In the point convention the same computation gives
$\sigma(E_1,P)=(E_1,[-p]P)=\langle-p\rangle'(E_1,P)$; the final answer below is unchanged because
$\chi_\ell$ is even and quadratic, but we fix the $\mu_M$-convention and use only the displayed formula.)

\emph{Galois action on the value.} $\overline G$ is defined over $\Fp$
(Proposition~\ref{prop:Gint}), so $\overline G(\sigma\tilde x)=\sigma\bigl(\overline G(\tilde x)\bigr)
=\overline G(\tilde x)^{p^2}$. Combining with the diamond identity of Proposition~\ref{prop:Gint},
\[
\overline G(\tilde x)^{p^2}=\overline G\bigl(\langle-p^{-1}\rangle\tilde x\bigr)
=\chi_\ell(-p^{-1})\,\overline G(\tilde x)=\chi_\ell(p)\,\overline G(\tilde x)=\leg{\ell}{p}\overline G(\tilde x),
\]
because $\chi_\ell(-1)=1$ and $\chi_\ell(p^{-1})=\chi_\ell(p)^{-1}=\chi_\ell(p)$. Since
$\overline G$ is invertible on $Y_1(M)_{\Fp}$, $\overline G(\tilde x)\ne0$ and
\begin{equation}\label{eq:Gfrob}
\overline G(\tilde x)^{p^2-1}=\leg{\ell}{p}.
\end{equation}
No coefficient-field factor enters: the only Galois action used is that of $\sigma$ on the point
$\tilde x$, and $\overline G$ is $\sigma$-invariant as a function.

\emph{Descent to $X_0(\ell)$.} By Proposition~\ref{prop:Gint},
$\overline G(\tilde x)^{24}=\overline D_\ell(\delta\tilde x)^{-1}=D_\ell(x)^{-1}$. Therefore
\[
D_\ell(x)^{k}=\overline G(\tilde x)^{-24k}=\overline G(\tilde x)^{-(p^2-1)}=\leg{\ell}{p}^{-1}=\leg{\ell}{p}.
\]
This holds at every geometric supersingular point, including those over $j=0$ and $j=1728$: the
argument takes place on the fine moduli scheme $Y_1(M)$, and $D_\ell(x)$ is defined by
\eqref{eq:Dmoduli} independently of automorphisms of $(E,C)$. For $\ell=N\in\mathcal N$,
$D_N=u_N^{N-1}$ and $(N-1)k=(p^2-1)/a_N$.
\end{proof}

\begin{remark}[comparison with Ono]\label{rem:ono}
\citet[Thm.~1.2]{Ono25} (v2, 3 July 2026) proves $R_\ell^{(p^2-1)/12}=1$ at supersingular points
for $R_\ell=\Delta(\ell\tau)/\Delta(\tau)=D_\ell^{-1}$, via $R_\ell(\tilde P)^{(p^2-1)/12}=\kappa(E')/\kappa(E)$
(his Lemma~4.4) and $\kappa(E)=\Delta(E,\omega)^{(p^2-1)/12}=(-1)^{(p+1)/2}$ on Deuring models with
$\Fq$-rational $\omega$ (his Thm.~1.3). His Remark~4 observes that the $(p^2-1)/24$ power can be $-1$
($p=11$, $\ell=2$: $R_2\in\{6,7,2\}$, fifth powers $-1$), that the values of
$g_\ell=\eta(\ell z)/\eta(z)$ need not lie in $\Fq$ (they generate $\mathbb F_{11^4}$ there), and that
a proof descending $\Delta^{1/24}$ or $\Delta^{1/12}$ to $\Fq$ must fail. Theorem~\ref{thm:C}
determines the sign: $R_\ell^{(p^2-1)/24}=\leg{\ell}{p}$, equivalently $R_\ell(\tilde P)$ is a
$24$-th power in $\Fq^\times$ iff $\leg{\ell}{p}=1$. This is consistent with Remark~4: $\leg{2}{11}=-1$,
and \eqref{eq:Gfrob} shows $\overline G(\tilde x)\in\mathbb F_{p^4}\setminus\Fq$ exactly when
$\leg{\ell}{p}=-1$. The proof above does not descend any root of $\Delta$ to $\Fq$: it uses that
$\overline G$ is an $\Fp$-rational \emph{function} on $Y_1(576\ell)$, and tracks its values through
the diamond character. The isogeny-invariance of $\kappa$ over $\Fq$ (for $E_1,E_2$ separably isogenous over $\Fq$,
the values $(\Lp{-1}\Delta(E_i))^{(p^2-1)/12}$ coincide) is already in \citet[\S2]{Baker98}, as a
consequence of the Gross--Landweber congruence; Ono's Theorem~1.2 follows from it by his Lemma~4.4.
\end{remark}

\section{The signed formula and a level-one corollary}\label{sec:proof}

\begin{proof}[Proof of Theorem~\ref{thm:main}(e)]
By Proposition~\ref{prop:product}, $r_iW_i=\beta(E)\,(\Delta/\Delta(N\tau))^k(x)$. Fact~\ref{fact:baker}
gives $\beta(E)=-\Lp{-1}$, and Theorem~\ref{thm:C} gives $(\Delta/\Delta(N\tau))^k(x)=\Lp N$. So
$r_iW_i=-\Lp{-1}\Lp N=-\Lp{-N}$. By Corollary~\ref{cor:rFp}, $r_i\in\Fp^\times$, and therefore
$W_i=-\Lp{-N}r_i^{-1}\in\Fp^\times$.
\end{proof}

\begin{proof}[Proof of Theorem~\ref{thm:main}(f)]
By Proposition~\ref{prop:res} and part (b), $\overline H(t^p,t)=-\Res_tR_{N,p}$. This is $0$ unless $t$ is thin, and equals
$-r_i=\Lp{-N}W_i^{-1}$ if $t=t_i$ is thin. This fixes the sign: with
$\Psi=(X^p-Y)(X-Y^p)+pH$, the residue is $-\overline H(t^p,t)$, not $+\overline H$.
\end{proof}

\begin{proof}[Proof of Corollary~\ref{cor:D}]
Combining Lemma~\ref{lem:Bisog} with Theorem~\ref{thm:C} gives, for every supersingular
$\ell$-isogeny $E\to E'$,
\[
\frac{c(E')}{c(E)}=\Lp\ell\,(\Delta/\Delta(\ell\tau))^k(x)=1.
\]
By Fact~\ref{fact:graph}, $c$ is constant. Fact~\ref{fact:baker} gives $c^2=\beta=-\Lp{-1}$.
Since $-\Lp{-1}$ is $-1$ for $p\equiv1\pmod4$ and $+1$ for $p\equiv3\pmod4$, $c\in\Fp$ in both cases.
For the congruence: the level-one mod-$p$ form $B^{(p-1)/2}-c\,\Delta^{(p^2-1)/24}$ vanishes on the supersingular locus.
The zeros of $A$ there are simple on the stack, including at $j=0,1728$. Hence $A$ divides it in the ring $\Fp[E_4,E_6]$ of level-one forms ($p\ge5$).
\end{proof}
\begin{remark}
Given Lemma~\ref{lem:Bisog}, Theorem~\ref{thm:C} for all primes $\ell$ is equivalent to the constancy of $c$. It therefore
follows from Theorem~\ref{thm:C} for a single $\ell$ together with Fact~\ref{fact:graph}.
\end{remark}

\begin{remark}[on Conjecture~\ref{conj:E}]
By Wilson's theorem $\bigl(\frac{p-1}2\bigr)!^2\equiv(-1)^{(p+1)/2}=-\Lp{-1}$, so the conjecture is
consistent with Corollary~\ref{cor:D}. Since $c(E)$ is constant, evaluation at
one supersingular curve for each $p$ would suffice. Choosing such a curve
uniformly is less immediate than this reduction suggests: $j=1728$ is
supersingular only for $p\equiv3\pmod4$, and $j=0$ only for
$p\equiv2\pmod3$. In particular neither familiar CM model is supersingular
when $p\equiv1\pmod{12}$. Moreover Baker's congruence fixes $c_p^2$ rather
than choosing between its two square roots; an explicit evaluation of $B$ is
still needed. The factorial formula holds in the recorded computations for all
primes $5\le p\le997$ and remains a conjecture here.
\end{remark}

\begin{remark}[$\Fp$-rationality at non-rational points]
The two rationality assertions in Theorem~\ref{thm:main}(e) have different explanations. $r_i\in\Fp$ is elementary (Corollary~\ref{cor:rFp}: symmetry
of $H$ and Frobenius-equivariance of residues). $W_i\in\Fp$ is not formal: writing
$W_i=\operatorname{Nm}_{\Fq/\Fp}(D_{\rm thin}'(t_i))\cdot\prod_{e_k>1}(t_i-t_k)^{(p+1)/e_k}$, it amounts to the product
$\prod_{e_k>1}(t_i-t_k)^{(p+1)/e_k}$ lying in $\Fp$. Equivalently, a product of power-residue characters (one of order $e_k$ for each factor) must be trivial; this is not visible factor by factor. Here it follows only
from part~(e).
\end{remark}


\section{Geometric interpretation of the weights}\label{sec:geometry}
The proof of Theorem~\ref{thm:main} is algebraic. We now interpret its weights in
the integral modular curve, retaining the geometric explanation of pole cancellation.
Put $W=W(\Fb)$ and let
\[
\mathcal B=X_0(Np)^{\mathrm{DR}}_W,\qquad
\mathcal A=\mathcal B/\langle w_N\rangle.
\]
Here $\mathcal B$ denotes the coarse Deligne--Rapoport model and $w_N$ is the
partial Atkin--Lehner involution. A split node with completed local ring
$W[[s,t]]/(st-p^ev)$, $v$ a unit, has \emph{thickness} $e$.

\begin{proposition}\label{prop:geometry}
The special fibre of $\mathcal A$ has two components, each isomorphic to
$X_0(N)^+_{\Fb}$. They meet at the supersingular Fricke orbits. At the orbit of
$x=(E,C_N)$ the node has thickness
\[
e_x=|\Aut(E,C_N)/\{\pm1\}|\,
|\Stab_{\langle w_N\rangle}(x)|.
\]
Consequently $R_{N,p}$ has a pole at a supersingular value precisely when the
corresponding node has thickness one.
\end{proposition}
\begin{proof}
Deligne--Rapoport's theorem applies with prime-to-$p$ level $N$ and the stabilizer
of a line in $(\Z/N)^2$ as the auxiliary group
\citep[V, 1.14--1.20; VI, Theorem~6.9, especially (6.9.1)]{DR73}.
It gives two copies of $X_0(N)_{\Fb}$ with Frobenius gluing, and the coarse
completed local ring at $x$ is
\[
W[[s,t]]/(st-p^{m_x}),\qquad m_x=|\Aut(E,C_N)/\{\pm1\}|.
\]
The forgetful map $\mathcal B\to X_0(N)_W$ is finite: it is proper and has finite
fibres on the generic fibre and on the two special-fibre components.
Normality identifies $\mathcal B$ with the normalization in its generic function
field. The generic partial Fricke map covers $w_N$ on $X_0(N)_W$, hence extends
uniquely by normalization. On the elliptic locus it is
\[
(E,C_N,C_p)\longmapsto(E/C_N,E[N]/C_N,q_N(C_p)).
\]
Since $q_N$ is prime to $p$, it preserves the Frobenius and Verschiebung components
and both branches at a node.

At a fixed node the stabilizer has order two and acts nontrivially on either
coarse branch. Indeed its action on the smooth curve $X_0(N)$ is a nontrivial
tame involution, whose derivative at a fixed point is $-1$.
Taking the branch-preserving quotient therefore multiplies the thickness by
this order, by \citet[Proposition~10.3.48(b)--(c), pp.~526--528]{Liu02}.
At a non-fixed node the inertia is trivial and the thickness is unchanged.
Invariants commute with reduction because $2$ is invertible in $W$.
The quotient components are consequently $X_0(N)^+_{\Fb}$, and no node is
smoothed, since no stabilizer exchanges branches. This proves the thickness formula.
The last assertion is Theorem~\ref{thm:main}(b).
\end{proof}

\begin{remark}[level two]
For $N=2$, the Fricke fixed points are $u=\pm64$, with $T=-104,152$.
At $u=-64$ the reduced automorphism group has order two and the Fricke stabilizer
has order two, giving thickness four. At $u=64$ their orders are one and two,
giving thickness two. All other supersingular values have thickness one.
The two exceptional values are supersingular precisely when
$\left(\frac{-1}{p}\right)=-1$ and $\left(\frac{-2}{p}\right)=-1$, respectively.
Thus a thick value is absent from the pole divisor but remains in the
weighted residue product. The same distinction applies at each of the five levels.
\end{remark}

\section{Finite certificates and computational evidence}\label{sec:comp}
The accompanying exact-arithmetic computations provide checks independent of the
signed product formula. They are not used in the proof. We describe the certification
bound so that the strength of each computation is explicit.

\begin{proposition}[finite certificate]\label{prop:certificate}
Let $L=p^2$, let $D\in\Fp[T]$ be a monic divisor of $T^L-T$, of degree $d$, and
let $P\in\Fp[T]$ have degree less than $d$. If
\[
D(\overline T)\,\overline{T|U_p}-P(\overline T)=O(q^{L-d+1}),
\]
then $R_{N,p}=P/D$. For $D=1$, take $P=0$.
\end{proposition}
\begin{proof}
Proposition~\ref{prop:res} implies that
$R_{N,p}-P/D=C/(T^L-T)$ with $\deg C<L$.
If $C\ne0$, its $q$-order is $L-\deg C$, since $T=q^{-1}+O(q)$.
Multiplication by $D(T)$ gives order $L-\deg C-d\le L-d$, contradicting the
hypothesis. Thus $C=0$.
\end{proof}

The first computation reconstructs $R_{N,p}$ from its $q$-expansion, using
$D=D_{\rm thin}$ built independently from supersingular data.
The coefficients of the numerator are recovered without the residue formula,
and Proposition~\ref{prop:certificate} then proves the reconstructed identity.
The reduced denominator was also recovered without a candidate for $p\le37$.
The Hecke polynomial was computed modulo $p^2$ for $p\le31$, checking symmetry,
the Kronecker congruence, and the residue criterion at every element of $\Fq$.

\begin{table}[ht]\centering
\begin{tabular}{rrrrr}\toprule
$N$ & certified pairs & thin values & non-rational thin values & range\\\midrule
2 & 45 & 530 & 310 & $5\le p\le211$\\
3 & 45 & 716 & 456 & $5\le p\le211$\\
5 & 44 & 1059 & 740 & $5\le p\le211$, $p\ne5$\\
7 & 44 & 1431 & 1116 & $5\le p\le211$, $p\ne7$\\
13 & 44 & 2510 & 2074 & $5\le p\le211$, $p\ne13$\\\bottomrule
\end{tabular}
\caption{Certified reconstructions. Thin values are distinct geometric values of $T_N$.}
\end{table}

All 222 admissible pairs with $p\le211$ passed the certificates, giving 6246 thin
values, of which 4696 are non-rational. Four pairs have no thin values; the ratio
set is empty in these cases. In every other pair it is a singleton equal to
$-\left(\frac{-N}{p}\right)$. A second calculation, using the double-pole formula
of Lemma~\ref{lem:g}, checks a further 960 pairs for $223\le p\le1499$.
That calculation depends on the derivative lemma and is not a separate
reconstruction from $T|U_p$.

The eta/Frobenius identity was also tested by a different route: normalized
V\'elu quotients and their discriminants, without eta functions or Hauptmoduln.
For the 15 primes $\ell\le47$ and all admissible $5\le p\le199$, the calculation
covers 647 pairs and 127541 cyclic-kernel checks on chosen representatives of
supersingular elliptic curves. These are not counts of distinct coarse points:
automorphisms may identify kernels over $j=0$ and $1728$.
Every kernel gave the asserted discriminant-ratio sign.

\paragraph{Data and software.}
The supplementary packages \texttt{Signed\_Weighted\_Residue\_v1\_0} and
\texttt{Eta\_Frobenius\_Lemma\_Audit\_v1\_0} contain Python code, per-pair JSON
results and reproduction instructions (\texttt{python-flint} 0.9.0 and
\texttt{mpmath}). The numerical ranges reported here are those recorded in these
packages. The present revision changes the proofs, attribution and presentation;
it does not assert a new full rerun of those ranges.

\section{Further directions}
The argument separates three inputs: the genus-zero coordinate and Kronecker
congruence, the Hasse divisor and Kodaira--Spencer derivative, and the signed
isogeny-ratio identity. Only the first two are specific to the pole theorem on the
five curves. Theorem~\ref{thm:C} applies to every prime isogeny degree different
from $p$, without a genus-zero assumption. Extending the residue theorem to other
Fricke curves requires an integral normalized coordinate and an explicit treatment
of the cusp divisor; it is not asserted here.

The constant $c_p$ of Corollary~\ref{cor:D} gives a level-one formulation of the
same sign information. Its factorial evaluation in Conjecture~\ref{conj:E}
remains open in this paper. The genus-zero moonshine identification supplies
examples of the modular functions to which the theorem applies; the result itself
is a statement about their reduction and supersingular divisors.

\paragraph{Acknowledgement of assistance.}
AI systems assisted in developing arguments, checking calculations and preparing
this manuscript. The source computations and the subsequent proof audit are
identified in the supplementary packages. This assistance does not constitute
independent specialist refereeing.

\begin{thebibliography}{99}\small
\bibitem[Apostol(1990)]{Apostol90} Apostol, T.~M. (1990). \emph{Modular Functions and Dirichlet Series in Number Theory}, 2nd ed. Graduate Texts in Mathematics 41. Springer.

\bibitem[Baker(1998)]{Baker98} Baker, A. (1998). A supersingular congruence for modular forms.
\emph{Acta Arithmetica}, 86(1), 91--100. \url{https://doi.org/10.4064/aa-86-1-91-100}.

\bibitem[Behrens(2009)]{Behrens09} Behrens, M. (2009). Congruences between modular forms given by the divided $\beta$ family in homotopy theory. \emph{Geometry \& Topology}, 13, 319--357. \url{https://doi.org/10.2140/gt.2009.13.319}.

\bibitem[Betina and Lecouturier(2018)]{BL18} Betina, A., \& Lecouturier, E. (2018). Congruence formulas for
Legendre modular polynomials. \emph{Journal of Number Theory}, 188, 71--87.

\bibitem[Chen et~al.(2019)]{CMT19} Chen, R.~C., Marks, S., \& Tyler, M. (2019).
$p$-Adic properties of Hauptmoduln with applications to moonshine.
\emph{SIGMA}, 15, 033, 35 pp. \url{https://doi.org/10.3842/SIGMA.2019.033}.

\bibitem[Conrad(n.d.)]{ConradQexp} Conrad, B. \emph{Algebraic theory of $q$-expansions}. Math 248B course handout, Stanford University.

\bibitem[Conway and Norton(1979)]{CN79} Conway, J.~H., \& Norton, S.~P. (1979). Monstrous moonshine.
\emph{Bulletin of the London Mathematical Society}, 11, 308--339.

\bibitem[Deligne and Rapoport(1973)]{DR73} Deligne, P., \& Rapoport, M. (1973). Les sch\'emas de modules de
courbes elliptiques. In \emph{Modular Functions of One Variable II} (LNM 349, pp.~143--316). Springer.

\bibitem[de~Shalit(1994)]{dS94} de~Shalit, E. (1994). Kronecker's polynomial, supersingular elliptic curves,
and $p$-adic periods of modular curves. \emph{Contemporary Mathematics}, 165, 135--148.

\bibitem[Deuring(1941)]{Deuring41} Deuring, M. (1941). Die Typen der Multiplikatorenringe elliptischer Funktionenk\"orper. \emph{Abhandlungen aus dem Mathematischen Seminar der Hansischen Universit\"at}, 14, 197--272.

\bibitem[Duncan and Swisher(2026)]{DS26} Duncan, J.~F.~R., \& Swisher, H. (2026).
Modular functions and the monstrous exponents. \emph{arXiv:2602.09135v1},
9 February 2026. \url{https://arxiv.org/abs/2602.09135v1}.

\bibitem[Gross(1990)]{Gross90} Gross, B.~H. (1990). A tameness criterion for Galois representations associated to modular forms (mod $p$). \emph{Duke Mathematical Journal}, 61(2), 445--517.

\bibitem[Jorgenson et~al.(2019)]{JST19} Jorgenson, J., Smajlovi\'c, L., \& Then, H. (2019).
The Hauptmodul at elliptic points of certain arithmetic groups.
\emph{Journal of Number Theory}, 204, 661--682.
\url{https://doi.org/10.1016/j.jnt.2019.03.021}.
The elliptic-value tables used here are also available in
\emph{arXiv:1602.07426v2}, \S4.2, Table~1 and Appendix~C.
\url{https://arxiv.org/abs/1602.07426v2}.

\bibitem[Kaneko and Zagier(1998)]{KZ98} Kaneko, M., \& Zagier, D. (1998). Supersingular $j$-invariants,
hypergeometric series, and Atkin's orthogonal polynomials. In \emph{Computational Perspectives on Number
Theory} (AMS/IP Studies in Advanced Mathematics 7, pp.~97--126).

\bibitem[Katz(1973)]{Katz73} Katz, N.~M. (1973). $p$-adic properties of modular schemes and modular forms.
In \emph{Modular Functions of One Variable III} (LNM 350, pp.~69--190). Springer.

\bibitem[Katz(1976)]{Katz76} Katz, N.~M. (1976). $p$-adic interpolation of real analytic Eisenstein series. \emph{Annals of Mathematics}, 104, 459--571.

\bibitem[Katz and Mazur(1985)]{KM85} Katz, N.~M., \& Mazur, B. (1985). \emph{Arithmetic Moduli of Elliptic
Curves} (Annals of Mathematics Studies 108). Princeton University Press.

\bibitem[Knopp(1970)]{Knopp70} Knopp, M.~I. (1970). \emph{Modular Functions in Analytic Number Theory}. Markham.

\bibitem[Koike(1973)]{Koike73} Koike, M. (1973). Congruences between modular forms
and functions and applications to the conjecture of Atkin.
\emph{Journal of the Faculty of Science, University of Tokyo, Section IA, Mathematics},
20, 129--169.

\bibitem[Landweber(1988)]{Landweber88} Landweber, P.~S. (1988). Supersingular elliptic curves and congruences for Legendre polynomials. In \emph{Elliptic Curves and Modular Forms in Algebraic Topology}, Lecture Notes in Mathematics 1326, pp.~69--93. Springer.

\bibitem[Larson(2020)]{Larson20} Larson, D.~M. (2020). Explicit modular forms from the divided beta family. \emph{Topology and its Applications}, 273, 107094. \url{https://doi.org/10.1016/j.topol.2020.107094}.

\bibitem[Liu(2002)]{Liu02} Liu, Q. (2002). \emph{Algebraic Geometry and Arithmetic Curves}. Oxford Graduate Texts in Mathematics 6. Oxford University Press.

\bibitem[Mestre(1986)]{Mestre86} Mestre, J.-F. (1986). La m\'ethode des graphes. Exemples et applications.
In \emph{Proceedings of the International Conference on Class Numbers and Fundamental Units of Algebraic
Number Fields} (pp.~217--242). Nagoya University.

\bibitem[Morton(2023)]{Morton23} Morton, P. (2023). Supersingular conjectures for the
Fricke group. \emph{International Journal of Number Theory}, 19(4), 803--828.
\url{https://doi.org/10.1142/S1793042123500409}; preprint
\url{https://arxiv.org/abs/2105.14532}.

\bibitem[Nakaya(2019)]{Nakaya19} Nakaya, T. (2019). The number of linear factors of
supersingular polynomials and sporadic simple groups. \emph{Journal of Number Theory},
204, 471--496. \url{https://doi.org/10.1016/j.jnt.2019.04.015}.

\bibitem[Ono(2004)]{Ono04} Ono, K. (2004). \emph{The Web of Modularity} (CBMS Regional Conference Series 102).
American Mathematical Society.

\bibitem[Ono(2026)]{Ono25} Ono, K. (2026). Modular forms for chromatic homotopy: supersingular congruences. \emph{arXiv:2509.16175v2}, 3 July 2026. \url{https://arxiv.org/abs/2509.16175v2}.

\bibitem[Pizer(1990)]{Pizer90} Pizer, A. (1990). Ramanujan graphs and Hecke operators. \emph{Bulletin of the
AMS (N.S.)}, 23(1), 127--137.

\bibitem[Robert(1980)]{Robert80} Robert, G. (1980). Congruences entre s\'eries d'Eisenstein, dans le cas supersingulier. \emph{Inventiones Mathematicae}, 61, 103--158.

\bibitem[Sakai(2011)]{Sakai11} Sakai, Y. (2011). The Atkin orthogonal polynomials for
the low-level Fricke groups and their application.
\emph{International Journal of Number Theory}, 7(6), 1637--1661.
\url{https://doi.org/10.1142/S1793042111004460}.

\bibitem[Sakai(2014)]{Sakai14} Sakai, Y. (2014). The Atkin orthogonal polynomials for
the Fricke groups of levels $5$ and $7$.
\emph{International Journal of Number Theory}, 10(8), 2243--2255.
\url{https://doi.org/10.1142/S1793042114500766}.

\bibitem[Waterhouse(1969)]{Waterhouse69} Waterhouse, W.~C. (1969). Abelian varieties over finite fields. \emph{Annales Scientifiques de l'\'Ecole Normale Sup\'erieure (4)}, 2, 521--560.
\end{thebibliography}
\end{document}
